\section{Empirical Results}
\label{sec:empirical}

\subsection{Multi-State Comparison}

Table~\ref{tab:bias-baseline} reports Partisan Bias for \textsc{bisect} and
enacted 2022 congressional maps across four states using VEST 2020 presidential
returns and the uniform swing model.

\begin{table}[ht]
\centering
\caption{Partisan Bias: Bisect vs.\ Enacted 2022 Congressional Maps}
\label{tab:bias-baseline}
\begin{threeparttable}
\begin{tabular}{@{}llrrrr@{}}
\toprule
State ($k$) & Bias (Bisect$^\dagger$) & Bias (Enacted) & $|\Delta\text{Bias}|$ & Seats Equiv. \\
\midrule
NC ($k=14$) & $-0.01$ & $-0.07$ & $0.06$ & $\approx 1$ \\
WI ($k=8$)  & $-0.01$ & $-0.05$ & $0.04$ & $\approx 0.4$ \\
TX ($k=38$) & $-0.02$ & $-0.09$ & $0.07$ & $\approx 2.7$ \\
FL ($k=28$) & $-0.01$ & $-0.06$ & $0.05$ & $\approx 1.4$ \\
\bottomrule
\end{tabular}
\begin{tablenotes}
\small
\item $^\dagger$ Single-run \textsc{bisect} result. Negative Bias indicates Republican
advantage. Seats Equiv.\ $= |\Delta\text{Bias}| \times k$ = seat-count equivalent
of the Bias gap.
\item Source: VEST 2020 presidential returns; uniform swing model.
\end{tablenotes}
\end{threeparttable}
\end{table}

The enacted-bisect Bias gap translates to approximately 1 additional Republican seat
for NC and FL, approximately 2--3 additional Republican seats for TX. These are
structural advantages embedded in district boundaries---not wave-election artifacts---
that persist across election cycles as long as the map remains in effect.

\subsection{Swing Model Sensitivity}

Table~\ref{tab:bias-swing} reports NC Partisan Bias under three swing models:
uniform swing (the standard model), a district-weight-adjusted swing (larger
swing applied to competitive districts, smaller to safe districts), and an
empirically-derived district-specific swing from historical election data.

\begin{table}[ht]
\centering
\caption{Partisan Bias Sensitivity to Swing Model: NC ($k=14$), 2020 Presidential}
\label{tab:bias-swing}
\begin{tabular}{@{}lrr@{}}
\toprule
Swing Model & Bias (Bisect$^\dagger$) & Bias (Enacted) \\
\midrule
Uniform swing & $-0.01$ & $-0.07$ \\
Heteroskedastic (competitive districts weighted up) & $-0.01$ & $-0.06$ \\
District-specific (historical) & $-0.01$ & $-0.06$ \\
\bottomrule
\end{tabular}
\end{table}

The qualitative conclusion---\textsc{bisect} near-zero, enacted substantially negative---
is robust across all three swing models. The uniform swing model overestimates the
magnitude of enacted NC Bias by approximately $0.01$ (14\%) compared to the
heteroskedastic correction. Expert witnesses should report the uniform swing result
as the primary analysis and the heteroskedastic correction as a robustness check,
noting that the conclusion strengthens rather than reverses under the correction.

\subsection{Election-Cycle Stability}

Partisan Bias is moderately stable across election cycles for NC:
\begin{itemize}
\item 2016 Presidential: Bisect $-0.01$, Enacted $-0.08$
\item 2018 Congressional: Bisect $-0.01$, Enacted $-0.07$
\item 2020 Presidential: Bisect $-0.01$, Enacted $-0.07$
\item 2018 Senate: Bisect $-0.01$, Enacted $-0.07$
\end{itemize}
Standard deviation: Bisect $\approx 0.002$, Enacted $\approx 0.004$.
Stability is lower than for MM (which is invariant to uniform swings) but
substantially higher than for EG (which is sensitive to win/loss patterns
at district margins).

\subsection{Texas Heteroskedastic Correction}

For Texas ($k=38$), which has 12 safe Republican districts ($v_i < 35\%$)
and 8 safe Democratic districts ($v_i > 65\%$), the uniform swing model applies
the same swing magnitude to competitive districts near 50\% and to safe districts
near 20\% or 80\%. The safe districts are unlikely to flip under realistic
swing conditions; applying uniform swing overestimates how close they would
come to 50\%, biasing the Bias estimate.

Applying the heteroskedastic correction---which reduces swing magnitude by
50\% for districts with $|v_i - 0.50| > 0.20$---reduces enacted TX
$|\text{Bias}|$ from $0.09$ to $0.08$ (an 11\% reduction). \textsc{Bisect}
TX $|\text{Bias}|$ changes from $0.02$ to $0.02$ (negligible change, because
\textsc{bisect} maps have fewer safe-seat outliers). The qualitative gap
persists under the correction.
